\documentclass[tikz,border=8pt]{standalone}
\usepackage{tikz}
\usepackage{amsmath}
\usepackage{amssymb}
\usepackage{pgfplots}
\pgfplotsset{compat=1.18}
\usepackage{ctex}
\usetikzlibrary{calc,positioning,arrows.meta}

% LC 42 Trapping Rain Water
% 双指针法：高度数组 [0,1,0,2,1,0,1,3,2,1,2,1]
% 蓝色柱子 + 阴影蓄水 + left/right 指针 + left_max/right_max 标注

\begin{document}
\begin{tikzpicture}

%% ── 标题 ──────────────────────────────────────────────────
\node[font=\large\bfseries] at (5.5, 4.6)
  {LC 42：接雨水（Trapping Rain Water）双指针法};

%% ── 参数 ─────────────────────────────────────────────────
% 高度数组 [0,1,0,2,1,0,1,3,2,1,2,1]，n=12
% 每格宽 0.9cm，高度单位 0.9cm
% heights: 0 1 0 2 1 0 1 3 2 1 2 1
% water:   0 0 1 0 1 2 1 0 0 1 0 0  => total = 6

%% ── 坐标系 ───────────────────────────────────────────────
\draw[->,gray!60,thin] (-0.3, 0) -- (11.2, 0) node[right,font=\scriptsize]{索引};
\draw[->,gray!60,thin] (0, -0.2) -- (0, 4.2)  node[above,font=\scriptsize]{高度};

% y 轴刻度
\foreach \h in {1,2,3}{
  \draw[gray!35,thin] (-0.1,\h*0.9) -- (10.9,\h*0.9);
  \node[font=\tiny,gray,left] at (-0.15,\h*0.9) {\h};
}

%% ── 柱子：index 0..11 ────────────────────────────────────
% height, water 分别手动列出，避免 \pgfmathparse 复杂性
% 格式：\drawbar{i}{height}{water}

% index 0: h=0, w=0
% index 1: h=1, w=0
\fill[blue!55] (0.85,0) rectangle (1.75,0.9);
\draw[blue!70!black,thin] (0.85,0) rectangle (1.75,0.9);
\node[font=\tiny,blue!80!black] at (1.35,1.05) {1};

% index 2: h=0, w=1
\fill[blue!20,opacity=0.85] (1.85,0) rectangle (2.65,0.9);
\draw[blue!40,thin,dashed] (1.85,0) rectangle (2.65,0.9);
\node[font=\tiny,blue!50!black] at (2.25,0.45) {+1};

% index 3: h=2, w=0
\fill[blue!55] (2.75,0) rectangle (3.55,1.8);
\draw[blue!70!black,thin] (2.75,0) rectangle (3.55,1.8);
\node[font=\tiny,blue!80!black] at (3.15,1.95) {2};

% index 4: h=1, w=1
\fill[blue!55] (3.65,0) rectangle (4.45,0.9);
\draw[blue!70!black,thin] (3.65,0) rectangle (4.45,0.9);
\fill[blue!20,opacity=0.85] (3.65,0.9) rectangle (4.45,1.8);
\draw[blue!40,thin,dashed] (3.65,0.9) rectangle (4.45,1.8);
\node[font=\tiny,blue!80!black] at (4.05,1.05) {1};
\node[font=\tiny,blue!50!black] at (4.05,1.5) {+1};

% index 5: h=0, w=2
\fill[blue!20,opacity=0.85] (4.55,0) rectangle (5.35,1.8);
\draw[blue!40,thin,dashed] (4.55,0) rectangle (5.35,1.8);
\node[font=\tiny,blue!50!black] at (4.95,0.9) {+2};

% index 6: h=1, w=1
\fill[blue!55] (5.45,0) rectangle (6.25,0.9);
\draw[blue!70!black,thin] (5.45,0) rectangle (6.25,0.9);
\fill[blue!20,opacity=0.85] (5.45,0.9) rectangle (6.25,1.8);
\draw[blue!40,thin,dashed] (5.45,0.9) rectangle (6.25,1.8);
\node[font=\tiny,blue!80!black] at (5.85,1.05) {1};
\node[font=\tiny,blue!50!black] at (5.85,1.5) {+1};

% index 7: h=3, w=0
\fill[blue!55] (6.35,0) rectangle (7.15,2.7);
\draw[blue!70!black,thin] (6.35,0) rectangle (7.15,2.7);
\node[font=\tiny,blue!80!black] at (6.75,2.85) {3};

% index 8: h=2, w=0
\fill[blue!55] (7.25,0) rectangle (8.05,1.8);
\draw[blue!70!black,thin] (7.25,0) rectangle (8.05,1.8);
\node[font=\tiny,blue!80!black] at (7.65,1.95) {2};

% index 9: h=1, w=1
\fill[blue!55] (8.15,0) rectangle (8.95,0.9);
\draw[blue!70!black,thin] (8.15,0) rectangle (8.95,0.9);
\fill[blue!20,opacity=0.85] (8.15,0.9) rectangle (8.95,1.8);
\draw[blue!40,thin,dashed] (8.15,0.9) rectangle (8.95,1.8);
\node[font=\tiny,blue!80!black] at (8.55,1.05) {1};
\node[font=\tiny,blue!50!black] at (8.55,1.5) {+1};

% index 10: h=2, w=0
\fill[blue!55] (9.05,0) rectangle (9.85,1.8);
\draw[blue!70!black,thin] (9.05,0) rectangle (9.85,1.8);
\node[font=\tiny,blue!80!black] at (9.45,1.95) {2};

% index 11: h=1, w=0
\fill[blue!55] (9.95,0) rectangle (10.75,0.9);
\draw[blue!70!black,thin] (9.95,0) rectangle (10.75,0.9);
\node[font=\tiny,blue!80!black] at (10.35,1.05) {1};

%% ── 下标标注 ─────────────────────────────────────────────
\foreach \i/\x in {
  0/0.5, 1/1.3, 2/2.25, 3/3.15, 4/4.05, 5/4.95,
  6/5.85, 7/6.75, 8/7.65, 9/8.55, 10/9.45, 11/10.35}{
  \node[font=\tiny,gray] at (\x,-0.28) {\i};
}

%% ── left/right 双指针标注（初始位置）────────────────────
% L=0 (蓝色↑), R=11 (红色↑)
\draw[->,>=Stealth,blue!80!black,thick,line width=1.2pt]
  (0.5,-0.65) -- (0.5,-0.05);
\node[font=\scriptsize,blue!80!black] at (0.5,-0.92) {L=0};

\draw[->,>=Stealth,red!80!black,thick,line width=1.2pt]
  (10.35,-0.65) -- (10.35,-0.05);
\node[font=\scriptsize,red!80!black] at (10.35,-0.92) {R=11};

%% ── left_max / right_max 标注 ─────────────────────────────
% left_max=2 (扫描左侧)
\draw[blue!60!black,dashed,thick]
  (-0.1,1.8) -- (3.65,1.8);
\node[draw=gray!50,fill=white,thin,rounded corners=2pt,
      font=\scriptsize,inner sep=2pt]
  at (1.5,2.12) {left\_max=2};

% right_max=2 (扫描右侧)
\draw[red!60!black,dashed,thick]
  (9.05,1.8) -- (11.0,1.8);
\node[draw=gray!50,fill=white,thin,rounded corners=2pt,
      font=\scriptsize,inner sep=2pt]
  at (9.9,2.12) {right\_max=2};

% 全局 max=3
\draw[green!60!black,dashed,thick]
  (5.0,2.7) -- (8.0,2.7);
\node[draw=gray!50,fill=white,thin,rounded corners=2pt,
      font=\scriptsize,inner sep=2pt]
  at (6.5,3.02) {peak=3 (i=7)};

%% ── 蓄水总量标注 ─────────────────────────────────────────
\node[draw=gray!50,fill=white,thin,rounded corners=2pt,
      font=\small\bfseries,inner sep=5pt]
  at (5.5,4.0) {蓄水总量 = 6 单位};

%% ── 算法说明框 ───────────────────────────────────────────
\node[draw=gray!50,fill=white,thin,rounded corners=2pt,
      font=\scriptsize,inner sep=5pt,align=left]
  at (5.5,-1.75) {%
    双指针：L=0，R=n-1；%
    若 h[L]$<$h[R]：water+=left\_max$-$h[L]，L++；%
    否则 water+=right\_max$-$h[R]，R--
  };

%% ── 图例 ─────────────────────────────────────────────────
\fill[blue!55]   (-0.05,-2.6) rectangle (0.35,-2.3);
\node[font=\scriptsize,right] at (0.4,-2.45) {柱子高度};
\fill[blue!20,opacity=0.85] (2.5,-2.6) rectangle (2.9,-2.3);
\draw[blue!40,thin,dashed] (2.5,-2.6) rectangle (2.9,-2.3);
\node[font=\scriptsize,right] at (2.95,-2.45) {蓄水区域};
\draw[->,>=Stealth,blue!80!black,thick] (5.5,-2.6) -- (5.5,-2.25);
\node[font=\scriptsize,right] at (5.6,-2.45) {left 指针 (L)};
\draw[->,>=Stealth,red!80!black,thick] (8.0,-2.6) -- (8.0,-2.25);
\node[font=\scriptsize,right] at (8.1,-2.45) {right 指针 (R)};

\end{tikzpicture}
\end{document}
